% Drafted against audit/outline.md §1. Every figure in this section is read
% from data/processed/diversification_comparison.csv,
% robustness_sign_variants.csv, robustness_windows.csv and
% robustness_age_weighting.csv in the frozen research repository; all four are
% hash-recorded in audit/evidence_freeze.csv at commit 1e6bdd74.

\section{Introduction}\label{sec:introduction}

Every large dealer that reports market risk under Item 7A of Form 10-K
publishes a table with roughly the same shape. Average daily Value-at-Risk is
broken out by risk category, the categories are added up, and a single negative
line is subtracted to arrive at the firmwide figure. Goldman Sachs calls that
line the diversification effect, and defines it as the difference between total
VaR and the sum of the VaRs of its four risk categories: interest rates, equity
prices, currency rates and commodity prices. Every other entry in the table is
a level. How much could be lost on rates, on equities, on currencies, on
commodities. The diversification line is not a level. It is the only number in
the table that depends on how those four exposures move against one another,
and co-movement is the one input the filing never publishes.

That asymmetry is what makes the line testable from outside, and it is also
what bounds any such test. A reader cannot check a disclosed interest-rate VaR
without knowing the rate book. A reader who builds a portfolio consistent with
the disclosed split of risk across categories, and exposes it to public market
data, can at least ask whether the co-movement that portfolio generates
resembles the co-movement the filing reports. That is the exercise here, and
the portfolio is a proxy throughout. It is calibrated so that each sleeve
carries the share of standalone risk that Goldman discloses for the
corresponding category; it is not a reconstruction of Goldman's positions, and
its dollar size is a normalisation that carries no information. What is being
tested is whether a static book with the disclosed risk structure, driven by
public factors, produces the disclosed diversification. Nothing in the paper
identifies a position, a desk, or a hedge.

The disclosures themselves make the test worth running, and they are more
specific about method than the VaR tables alone suggest.
\citet{jorion2002informative} showed that bank VaR disclosures are informative
about subsequent trading volatility, which is the premise any replication
presumes. The literature that followed has concentrated on the quality and the
conservatism of the reported numbers: \citet{perignon2010level} graded
disclosure practice across a cross-section of commercial banks,
\citet{perignon2008overstate} asked whether the levels banks report are
systematically too high, and \citet{perignon2010diversification} examined the
diversification component directly, on quarterly observations for four US banks
running to the first quarter of 2007. Those studies work across institutions,
and at the reporting frequency their data allow.

This paper adds three things to that base. The first is a single-dealer series
long enough to contain regime breaks rather than straddle one. Goldman's
diversification effect is reconstructed from ten consecutive filings covering
fiscal years 2016 through 2025, a decade that includes the March 2020
dislocation and the 2022--2023 tightening cycle. The second is a replication
run at the bank's own disclosed settings. Every filing in the sample states
that Goldman computes VaR by historical simulation with full valuation, samples
its scenarios from five years of historical data, and weights that data so that
the importance of older observations decays. All three are implemented here, on
seventeen factors drawn from the US Treasury and ECB curves and from
exchange-traded prices, with the interest-rate sleeve valued by full bond
revaluation rather than by a duration approximation. The engine is therefore
capable of disagreeing with the disclosure in a way that an identity built out
of the disclosed numbers cannot, and when it does disagree the disagreement
cannot be attributed to a cruder method. The third is a decomposition. Goldman discloses the same quantity twice, as an
annual average and as a period-end snapshot, and comparing each against its
proxy analogue separates how far the ratio moves from whether those movements
persist. The separation is what turns a failure to reproduce stability into a
specific claim about mean reversion, and it uses only what the filings already
publish. Running the complete space of sign configurations for the four sleeves
then bounds how closely a static book can match the level and the steadiness at
once, which turns out to be the binding constraint.

The disclosed series is remarkably flat. Across the ten years the
diversification effect ranges from 39.0\% to 45.5\% of summed category VaR,
with a mean of 42.0\% and a standard deviation of 2.08 percentage points. A
regression on the year index finds no trend ($p = 0.23$). A decade containing a
pandemic, a zero-rate period and more than 500 basis points of policy
tightening moved the reported co-movement of Goldman's four risk categories by
about two points of standard deviation, and moved it in no particular
direction.

The replication reproduces that level and does not reproduce that steadiness.
At the disclosed five-year window, a book sized once on 2019 and then left
alone averages 41.6\% against the disclosed 42.0\%; the two levels are
statistically indistinguishable ($p = 0.84$). Its standard deviation across
those years is 6.33 percentage points against the disclosed 2.08, a difference
a variance-ratio test puts at $p = 0.0028$. And the two move differently over
the decade: Goldman's second-half mean sits 0.8 points above its first, while
the proxy's sits 8.5 points below.

Goldman's second disclosed measure of the same quantity says what kind of
failure it is. Alongside the average daily figure the filings give a period-end
snapshot, and that series has a standard deviation of 6.72 points: the firm's
diversification ratio moves a great deal, and what is steady is its mean over a
year rather than the ratio itself. Measured period-end against period-end the
proxy's variation is statistically indistinguishable from Goldman's. What
separates them is persistence. The lag-1 autocorrelation of consecutive annual
values is $-0.47$ for Goldman's average daily series and essentially zero for
its period-end series, against roughly $+0.9$ for both proxy series: Goldman's
movements do not carry forward and the proxy's do. What fails to transfer is
therefore not stability but mean reversion, and mean reversion is what the
disclosed steadiness consists of.

The natural objection is that a proxy book on seventeen factors is simply a
blunter instrument than the bank's, and that the gap measures the crudeness of
the replication. Adopting Goldman's disclosed age weighting tests that
objection directly, and refutes it. Weighting the scenario set toward recent
observations raises the year-to-year standard deviation from 6.33 to between
9.46 and 10.04 points, four to five times the disclosed figure, across decay
rates with half-lives of three weeks and three months. It does so while making
the engine a demonstrably better risk model: the exception rate falls from
6.98\% to 5.45\% against a nominal 5\%, and the Kupiec proportion-of-failures
test moves from rejecting the model below any conventional level to a $p$-value
of 0.328. The disclosed weighting is the correct specification, the filings say
so, and applying it widens the stability gap rather than closing it.

Nor is the gap an artefact of one calibration of the sleeves. The four can be
signed in sixteen ways and all sixteen were run. Taking the disclosed standard
deviation itself as the bar, five configurations are at least that steady and
the highest of those averages 35.64\%, six points below the disclosure; eight
average 40\% or more and the steadiest of those varies by 3.68 points a year.
No configuration achieves both. Within this engine the two properties of the
disclosed series trade off against each other, and that frontier, rather than
any single calibration, is the paper's central result. It also points at what a
static book omits. A dealer that reprices positions daily, and whose hedges
respond to the correlations it observes, can hold a diversification ratio
steady in a way that a fixed portfolio exposed to the same market cannot.
Goldman's own filings attribute its year-over-year VaR changes to exposure
movements as well as to volatility in four consecutive years.

The remainder of the paper is organised as follows. Section~\ref{sec:disclosed}
sets out what Goldman discloses, how the series was extracted from the filings,
and where this exercise sits relative to prior work.
Section~\ref{sec:sample} describes the factor panel and the choices that shaped
it. Section~\ref{sec:method} specifies the historical-simulation engine, the
disclosed settings it adopts, and the calibration that ties the proxy book to
the disclosed split. Section~\ref{sec:validation} reports the validation:
unit-level checks, independent reproduction of the backtest statistics, and the
formal coverage tests applied to the engine's own output.
Section~\ref{sec:results} presents the comparison against the disclosures and
Section~\ref{sec:assumptions} isolates which assumptions the stability result
depends on. Section~\ref{sec:limitations} states what the exercise cannot
establish, and Section~\ref{sec:conclusion} concludes. Code, extracted
disclosures and model output are available in the repository cited in the data
availability statement.
